\begin{ouverture}[calcul]{8}{Pourcentages et grandeurs quotients}{Partie 2 · Algèbre et modèles}
{Coefficient multiplicateur · Vitesse · Débit · Masse volumique}
\objectif{Calculer avec des pourcentages et un coefficient multiplicateur.}
\objectif{Calculer et interpréter une grandeur quotient avec son unité.}
\objectif{Convertir des durées, des vitesses et des débits.}
\prerequis{Proportionnalité et rapports (chapitre 7), conversions usuelles.}
\end{ouverture}
\pagestyle{otMain}\essentiel\fontsize{11}{15}\selectfont
\phantomsection\label{cours-8-1}
\essSection{1}{Pourcentages et coefficient multiplicateur}
\begin{propriete}Prendre $t\,\%$ d’une quantité revient à la multiplier par $\frac t{100}$. Pour exprimer une part en pourcentage, on divise la partie par le total : 28 élèves sur 80, c’est $\frac{28}{80}=0{,}35=35\,\%$.\end{propriete}
\begin{definition}Augmenter une grandeur de $t\,\%$ revient à la multiplier par $1+\frac t{100}$ ; la diminuer de $t\,\%$ revient à la multiplier par $1-\frac t{100}$. Ce nombre est le coefficient multiplicateur : $1{,}15$ pour une hausse de 15\,\%, $0{,}8$ pour une baisse de 20\,\%.\end{definition}
\begin{exemple}Un article de 60 € baisse de 15\,\% : il coûte $60\times0{,}85=51$ €. La baisse vaut $60-51=9$ €, soit bien $\frac{15}{100}\times60$. Retrouver le prix initial à partir de 51 € revient à diviser par 0,85.\end{exemple}
\begin{attention}Pour deux évolutions successives, on multiplie les coefficients : une baisse de 20\,\% puis de 10\,\% donne $0{,}8\times0{,}9=0{,}72$, soit une baisse de 28\,\% et non de 30\,\%. Pour réunir deux groupes, on revient aux effectifs avant de calculer un pourcentage.\end{attention}
\clearpage
\phantomsection\label{cours-8-2}
\essSection{2}{Grandeurs quotients}
\begin{definition}Une grandeur quotient est le quotient de deux grandeurs ; son unité combine les deux unités. Vitesse moyenne : $v=\frac dt$, en km/h ou m/s. Débit : $q=\frac Vt$, en L/min. Masse volumique : $\rho=\frac mV$, en g/cm$^3$. Prix au kilogramme : en €/kg.\end{definition}
\begin{exemple}Parcourir 150 km en 2 h 30 min donne $t=2{,}5$ h, puis $v=\frac{150}{2{,}5}=60$ km/h. Un robinet qui fournit 12 L en 40 s a un débit de $\frac{12}{40}=0{,}3$ L/s.\end{exemple}
\begin{savoirfaire}[Choisir des unités cohérentes]$1$ h $=60$ min $=3\,600$ s ; 1 h 30 min vaut 1,5 h et non 1,30 h. $1$ km/h $=\frac{1\,000}{3\,600}$ m/s : pour passer de km/h à m/s, on divise par 3,6. $1$ L $=1$ dm$^3$ et $1$ m$^3=1\,000$ L. On convertit séparément les deux unités d’une grandeur quotient.\end{savoirfaire}
\begin{attention}La vitesse moyenne d’un aller-retour n’est pas en général la moyenne des deux vitesses : on divise la distance totale par la durée totale. Écrire les unités à côté des données permet de contrôler un résultat.\end{attention}

\clearpage\rubrique{Découverte}
\begin{blocLivre}{ex-8-1}\begin{exercice}[1]{Une part en pourcentage}
Calculer 18\,\% de 250.
\end{exercice}\end{blocLivre}
\begin{blocLivre}{ex-8-2}\begin{exercice}[1]{Retrouver le pourcentage}
Dans un groupe de 32 élèves, 12 sont demi-pensionnaires. Quel pourcentage cela représente-t-il ?
\end{exercice}\end{blocLivre}
\begin{blocLivre}{ex-8-3}\begin{exercice}[1]{Des coefficients}
Donner le coefficient multiplicateur associé à une hausse de 15\,\%, une baisse de 20\,\%, une hausse de 3\,\% et une baisse de 7,5\,\%.
\end{exercice}\end{blocLivre}
\begin{blocLivre}{ex-8-4}\begin{exercice}[1]{Une vitesse moyenne}
Un cycliste parcourt 45 km en 2 h 30 min. Calculer sa vitesse moyenne.
\end{exercice}\end{blocLivre}
\begin{blocLivre}{ex-8-5}\begin{exercice}[1]{Une durée}
Convertir 2 h 18 min en heures décimales et en minutes.
\end{exercice}\end{blocLivre}

\rubrique{Entraînement}
\begin{blocLivre}{ex-8-6}\begin{exercice}[2]{Une remise}
Un sac coûte 80 euros. On lui applique une remise de 12\,\%. Calculer son nouveau prix avec un coefficient multiplicateur.
\end{exercice}\end{blocLivre}
\begin{blocLivre}{ex-8-7}\begin{exercice}[2]{Une hausse de loyer}
Un loyer de 640 euros augmente de 2,5\,\%. Calculer le nouveau loyer avec un coefficient multiplicateur, puis le montant de la hausse.
\end{exercice}\end{blocLivre}
\begin{blocLivre}{ex-8-8}\begin{exercice}[2]{Une conversion de vitesse}
Convertir 72 km/h en m/s.
\end{exercice}\end{blocLivre}
\begin{blocLivre}{ex-8-9}\begin{exercice}[2]{Changer les deux unités}
Convertir un débit de $0{,}012$ m$^3$/s en L/min.
\end{exercice}\end{blocLivre}
\begin{blocLivre}{ex-8-10}\begin{exercice}[2]{Une masse volumique}
Une pièce a un volume de 40 cm$^3$ et une masse de 108 g. Calculer sa masse volumique, puis la masse d’une pièce de 65 cm$^3$ du même matériau.
\end{exercice}\end{blocLivre}

\rubrique{Problèmes}
\begin{blocLivre}{ex-8-11}\begin{exercice}[3]{Deux remises successives}
Un article de 100 euros subit une baisse de 20\,\%, puis une baisse de 10\,\% du nouveau prix. Quel prix obtient-on ? La baisse totale vaut-elle 30\,\% ?
\end{exercice}\end{blocLivre}
\begin{blocLivre}{ex-8-12}\begin{exercice}[3]{Un pourcentage inverse}
Après une remise de 25\,\%, un objet coûte 54 euros. Retrouver son prix avant remise.
\end{exercice}\end{blocLivre}
\begin{blocLivre}{ex-8-13}\begin{exercice}[3]{Deux groupes réunis}
Un club compte 20 jeunes dont 40\,\% de filles, et 30 adultes dont 60\,\% de femmes. Quel est le pourcentage de personnes de sexe féminin dans le club ?
\end{exercice}\end{blocLivre}
\begin{blocLivre}{ex-8-14}\begin{exercice}[3]{Aller et retour}
Un trajet de 30 km est parcouru à 15 km/h à l’aller et à 30 km/h au retour. Calculer la vitesse moyenne sur l’ensemble du trajet.
\end{exercice}\end{blocLivre}
\begin{blocLivre}{ex-8-15}\begin{exercice}[3]{Une comparaison de débits}
La pompe A débite 900 L/min ; la pompe B débite $0{,}018$ m$^3$/s. Laquelle est la plus rapide ? Combien de minutes gagne-t-elle pour pomper 5 400 L ?
\end{exercice}\end{blocLivre}

\clearpage
